\documentclass[11pt,a4paper]{article}
\usepackage[T1]{fontenc}
\usepackage{lmodern}   % Type 1 vector fonts (replaces PK bitmaps)
\usepackage{cmap}      % ToUnicode CMap -> copyable text and accents
\usepackage[utf8]{inputenc}
\usepackage[margin=2.3cm]{geometry}
\usepackage{amsmath,amssymb}
\usepackage{booktabs,tabularx}
\usepackage{enumitem}
\usepackage[hidelinks]{hyperref}
\setlist{itemsep=2pt,topsep=3pt}
\newcolumntype{Y}{>{\raggedright\arraybackslash}X}

\title{Short paper: literature report\\[2pt]\large Cartels, leniency and damages in Brazil --- the three papers in \texttt{Artigos/}}
\author{Economia das Organizações e dos Incentivos --- MPE Insper 2026}
\date{24 November 2026 (Lecture 8: three papers and the idea of the model)}

\begin{document}
\maketitle

\section*{Purpose}
This report prepares the Lecture 8 deliverable of the short paper (guide:
\url{Aula/Short_paper_guia_rubrica.pptx}). For each paper the guide asks for five things:
the paper's question, its link to game theory, its mechanism or model, its main result, and how it can be
adapted to the group's research question. The short paper is theoretical: empirical results enter
only through the motivation and the literature review. The model, its equilibrium and its
comparative statics must be the group's own work. Section~\ref{sec:model} therefore offers
directions, not a finished model.

\begin{table}[h]
\centering\small
\begin{tabularx}{\textwidth}{@{}lYp{2.1cm}p{2.3cm}@{}}
\toprule
 & Paper & Type & Role \\
\midrule
1 & Sabino, Lara \& Sarto (2024), ``Game Theory Applied to the Brazilian Antitrust Leniency Program'', \emph{RDP} 21(111), 330--357, doi:10.11117/rdp.v21i111.7589 & Theory & Precedent and gap \\
2 & Fraga, Zenkner \& Micheli (2023), ``Responsabilização 360°'', \emph{Rev.\ Defesa da Concorrência} 11(1), 169--185, doi:10.52896/rdc.v11i1.1035 & Legal, descriptive & Institutions, parameters \\
3 & Oliveira \& Lima (2025), ``Estimativa de Sobrepreço em Cartéis: Evidências do Cartel do Cimento no Brasil'', CADE/DEE Working Paper 003 & Empirical (DiD) & Motivation, evidence \\
\bottomrule
\end{tabularx}
\end{table}

\noindent\textbf{Common thread.} A cartel is a repeated game that firms play against each other and
against an enforcer. Fines, leniency, damages and debarment all work by changing its payoffs.

%-------------------------------------------------------------
\section{Paper 1 --- the Brazilian leniency game}

\paragraph{Question.} Do the benefits of CADE's leniency programme (Law 12,529/2011, arts.\ 86--87)
lead a cartel member to propose an agreement? Would immunity from civil damages change that?

\paragraph{Link to the course.} A sequential game with perfect information, Nash equilibrium versus
subgame-perfect equilibrium, and backward induction, with Nature replaced by expected payoffs
(Lectures 3--4; Tadelis Ch.~8).

\paragraph{Model (pp.\ 344--353).} The Applicant chooses Propose or Not propose. After a proposal,
CADE chooses Accept or Reject. Detection probability $p=0.5$, an assumption the authors state
openly. Values in millions, for a firm with revenue 100:
cartel profit $L=20$, fine $M=20$, administrative sanction $A=5$, litigation cost $C=1$,
investigation cost $I=0.1M=2$, fines on the four other members $M_i=4L=80$, damages $RC=L=20$
(doubled to $2RC$ without an agreement, under Law 14,470/2022).
\begin{align*}
X_{\text{NP}} &= p(L-A-M-2RC-C)+(1-p)L = 0.5(-46)+0.5(20) = -13,\\
Y_{\text{NP}} &= p(A+M_i+M-I)+(1-p)(-L) = 0.5(103)-10 = 41.5.
\end{align*}
After Propose and Accept, CADE gets $80$ and the Applicant gets $-RC=-20$ (Game 1), or $0$ with
damages immunity (Game 2). After a rejection, both get the no-proposal payoffs.

\paragraph{Result.} Game 1: the SPE is (Not propose, Accept), so current law does not induce
leniency. Game 2: the Nash equilibria are (Propose, Accept) and (Not propose, Reject), and the SPE
is (Propose, Accept). The authors recommend damages immunity for applicants and warn that Law
14,470/2022, by encouraging damages claims, may weaken the incentive to report.

\paragraph{Re-solution (ours, to be checked by the group).} The cartel profit $L$ enters both
no-proposal branches but is dropped from the leniency branch. Counted consistently, the leniency
payoff is $L-RC=0$ under current law and $L=20$ with immunity. Since $X_{\text{NP}}(p)=20-66p$,
the Applicant proposes iff its leniency payoff is at least $20-66p$:

\begin{center}\small
\begin{tabular}{@{}lcc@{}}
\toprule
Leniency payoff & Propose iff & At $p=0.5$ \\
\midrule
Paper, Game 1: $-20$ & $p\ge 0.606$ & No \\
Paper, Game 2: $0$ & $p\ge 0.303$ & Yes \\
Consistent, Game 1: $L-RC=0$ & $p\ge 0.303$ & \textbf{Yes (result reverses)} \\
Consistent, Game 2: $L=20$ & always (weakly) & Yes \\
\bottomrule
\end{tabular}
\end{center}

\paragraph{Limitations.}
\begin{itemize}
\item No interaction among cartel members: there is no race to report, and $p$ is fixed from outside.
\item CADE's move is trivial, since Accept dominates, and its payoff is fines collected, not welfare.
\item No cartel-formation stage, although the paper draws conclusions about deterrence.
\item No comparative statics.
\item Minor slips: $A$ is described as a quarter of revenue (25) but set to 5; ``RC = M'' is written
for $RC=L$; ``40.5'' is written for 41.5.
\end{itemize}

\paragraph{Use.} The closest theoretical precedent. Its limitations define the gap.

%-------------------------------------------------------------
\section{Paper 2 --- ``Responsabilização 360°''}

\paragraph{Question.} Does CADE push for cartel liability beyond antitrust, in the criminal, other
administrative and civil spheres?

\paragraph{Method and result (pp.\ 173--174).} A search of CADE's case system, January 2018 to
April 2023. Of 40 bid-rigging cases judged, with R\$~2.52~bn in fines, 24 (60\%) referred the
decision to other bodies: prosecutors, audit courts, the Federal Police, victims. Of these cases, 26
concerned public procurement, and 20 of those ended in conviction. In 80\% of the 20 convictions
CADE engaged in ``360° accountability''; they total R\$~2.3~bn in fines. The precedents described
include Claro, Oi and Telefônica (about R\$~800~m, with a referral to the Correios), Infraero
cafeterias (five-year debarment) and Juazeiro do Norte (a civil servant involved). The count is
descriptive and makes no causal claim.

\paragraph{Institutional parameters (pp.\ 179--182).} The same conduct is an offence under the
Anti-Corruption Law (12,846/2013), a crime under the old procurement law (art.\ 90, two to four
years), and improbity when an official takes part (8,429/1992). Law 14,470/2022 makes five changes:
\begin{itemize}
\item damages claims expire five years after CADE's decision;
\item pass-on is not presumed;
\item victims get double damages, except against signatories of agreements with CADE, who pay single damages;
\item signatories are not jointly liable for the other members' damage;
\item CADE's decision is enough to ground a fast-track ruling (\emph{tutela de evidência}).
\end{itemize}

\paragraph{Link to the course.} The paper is indirect on theory. More sanction spheres and
debarment raise the expected sanction, which tightens the incentive constraint for collusion.
They also raise the cost of confessing, unless leniency extends to those spheres. This is the
tension behind Paper~1.

\paragraph{Use.} Institutional evidence and model parameters. It should not be presented as one of
the theory papers.

%-------------------------------------------------------------
\section{Paper 3 --- the cement cartel overcharge}

\paragraph{Question.} What overcharge did the cement cartel impose, and what did consumers save
when it broke? The cartel was convicted in May 2014 with R\$~3.11~bn in fines, the largest CADE
has imposed, and was active since at least 1987.

\paragraph{Method (pp.\ 18--25).} Difference-in-differences with an event study.
\[
\ln P_{ist} = \alpha + \beta\,(\text{Cement}_i\times\text{Post}_t) + \gamma_i + \mu_s + \tau_t + \varepsilon_{ist}.
\]
\begin{itemize}
\item Treated product: cement CP~II-E-32 per kg. Controls: CA-50 and CA-60 steel bars and hydrated
lime (same unit, no input link, not cartelised).
\item Data: monthly median prices by state from SINAPI-IBGE.
\item Break: the February 2007 dawn raids. The window is kept short to exclude the PAC demand boom
(January 2007).
\end{itemize}

\paragraph{Results (pp.\ 26--46).}
\begin{itemize}
\item There are no significant pre-trends; the effect becomes significant from April 2007.
\item Preferred estimate (June 2006 to August 2007): $\hat\beta=-0.0204$, so
$e^{\hat\beta}-1=-2.02\%$. The range across specifications is 1.40--2.35\%.
\item Robustness: 1.55--2.74\% with other control groups, and 1.58--1.72\% with CBIC data.
\item Estimated consumer savings, 2007--2022: about R\$~11.8~bn. In alternative windows with 2.02,
10 and 20\% overcharges: R\$~2.13/10.55/21.15~bn for 2007--2010, 5.84/28.90/57.80~bn for
February 2007 to May 2014, and 8.95/44.33/88.66~bn to May 2019.
\end{itemize}
The authors read 2.02\% as a lower bound. A meta-analysis gives a mean overcharge of 15.47\% and
a median of 16.01\% (Boyer \& Kotchoni 2015), and a German cement cartel shows 20.3--26.2\%.

\paragraph{Link to the course (strong).} The paper explains the low estimate with repeated-game
theory (Lecture 5):
\begin{itemize}
\item cartels price below monopoly to avoid detection (Harrington 2004b, 2005);
\item firms keep post-cartel prices high to shrink the damages counterfactual (Harrington 2004a);
\item collusion continues tacitly after the cartel ends (prices rise from September 2007; Turner
2024 finds the same in air cargo);
\item price wars arise in equilibrium (Green \& Porter 1984; Rotemberg \& Saloner 1986).
\end{itemize}
The cartel was exposed by an employee whistleblower, not by corporate leniency.

\paragraph{Use.} Motivation, with a real case and numbers, and evidence that damages rules feed
back into pricing.

%-------------------------------------------------------------
\section{Theory behind the three papers}
\begin{itemize}
\item \textbf{Expected sanction} (Becker 1968): deterrence compares the gain with $p\times$sanction.
Paper~1 models this directly. Paper~2 enlarges the sanction. Paper~3 measures the gain.
\item \textbf{Collusion in a repeated game} (Tadelis Ch.~10): with $n$ Bertrand firms and grim
trigger strategies, collusion is sustainable iff $\delta\ge 1-1/n$. In general,
$\delta \ge (\pi^{D}-\pi^{C})/(\pi^{D}-\pi^{N})$.
\item \textbf{Enforcement} (illustration): subtracting an expected sanction $pK$ per period gives
$\delta \ge 1-1/n + pK/\pi^{M}$, so detection and sanctions raise the bar for collusion.
\item \textbf{Leniency}: the first reporter can cheat and report at once, which raises the payoff
from deviating and creates a prisoner's-dilemma race (Motta \& Polo 2003; Spagnolo 2004; Aubert,
Rey \& Kovacic 2006; Harrington 2008). Reduced fines can, however, also lower expected sanctions
and stabilise the cartels that survive (Bigoni et al.\ 2012).
\item \textbf{Damages versus leniency}: damages deter collusion, but damages on the applicant lower
the value of confessing (Marvão \& Spagnolo 2018; Buccirossi, Marvão \& Spagnolo 2020). Law
14,470/2022 widens the gap between reporting and being reported on.
\end{itemize}

%-------------------------------------------------------------
\section{Research question and model directions}\label{sec:model}

\paragraph{Assessment of the set.} The set has one weak theory paper, one institutional paper and
one empirical paper. The literature criterion (25\% of content) rewards separating theory from
evidence and linking the theory to the model, so we recommend adding (or replacing Paper~2 with)
a canonical theory paper. Candidates already cited inside Papers 1 and 3 are Buccirossi, Marvão
\& Spagnolo (2020), Pinha \& Braga (2021, Brazil-specific) and Motta \& Polo (2003). Each must be
checked by DOI before it is cited.

\paragraph{Candidate question.} \emph{Under Law 14,470/2022 (double damages for other members,
single damages for the leniency applicant), does private damages exposure strengthen or weaken the
race to report, and thus make cartels less or more stable?}

\paragraph{Directions (the group chooses and solves).}
\begin{enumerate}[label=\Alph*.]
\item \emph{Static reporting game} (Lectures 1--2): two members choose Report or Silent at the
same time, with $p$, a fine $F$, a damages multiplier $k\in\{1,2\}$ and a leniency discount.
When is reporting dominant?
\item \emph{Repeated collusion with leniency} (Lectures 5--6): solve for
$\delta^*(p,F,k,\text{debarment})$ with and without a reporting option, and do comparative statics
on the damages gap. This is the most complete option for the model criterion (30\%).
\item \emph{Sequential} (Lecture 4): Paper~1 corrected, with a second member who can report first,
so that $p$ is no longer fixed.
\end{enumerate}

%-------------------------------------------------------------
\section{Further reading}
\small
Cited inside the three papers: Motta \& Polo (2003, \emph{IJIO} 21(3), 347--379); Buccirossi,
Marvão \& Spagnolo (2020, \emph{J.\ Legal Stud.} 49(2)); Pinha \& Braga (2021, \emph{J.\ Comp.\ Law
\& Econ.} 17(4), 877--902); Marvão \& Spagnolo (2018, \emph{Handbook of Game Theory and IO} vol.~2);
Bigoni et al.\ (2012, \emph{RAND} 43(2)); Bodnar et al.\ (2023, \emph{JLEO} 39(1)); Leslie (2006,
\emph{J.\ Corp.\ Law} 31); Harrington (2004a, ``Post-cartel pricing during litigation'', \emph{JIE});
Harrington (2004b, \emph{RAND}; 2005, \emph{IER}); Green \& Porter (1984, \emph{Econometrica});
Rotemberg \& Saloner (1986, \emph{AER}); Levenstein \& Suslow (2006, \emph{JEL}); Boyer \& Kotchoni
(2015, \emph{RIO}).

\smallskip\noindent Standard references not cited in the papers (verify before use): Becker (1968,
\emph{JPE}); Spagnolo (2004, CEPR DP ``Divide et Impera''); Aubert, Rey \& Kovacic (2006,
\emph{IJIO} 24); Harrington (2008, ``Optimal corporate leniency programs'', \emph{JIE} 56);
Porter (1983, \emph{Bell J.\ Econ.}). Textbooks: Tadelis (2013) Ch.~8 (pp.\ 151--174) and Ch.~10
(pp.\ 190--219); Cabral, \emph{Introduction to Industrial Organization}, \S7.3.

\sloppy\medskip\noindent\textbf{AI declaration note.} This report and the companion audio script
(\url{Leituras/short-paper-articles-narrated.txt}) were prepared with AI assistance for reading
and summarising. The guide requires the group to verify every reference and every calculation,
including the re-solution in Section~1, and to declare this use in the final slide.

\end{document}
